O modelo de Média Móvel Integrada Auto-Regressiva com Sazonalidade (SARIMA) é uma extensão do modelo de Média Móvel
Integrada Auto-Regressiva (ARIMA) que combina três elementos específicos em sua composição, a autoregressão (AR), a
média móvel (MA) e a diferenciação dos valores consecutivos da série (I), sendoexpandido para lidar com padrões sazonais
\cite{box}.

A componente de autoregressão descreve a dependência de uma observação atual e observações anteriores, em defasagens
$p$. Em um modelo puramente autorregressivo, uma nova observação $X_t$ é expressa como uma combinação linear dos $p$
valores anteriores, ponderados por coeficientes $\alpha_i$, conforme a equação \autoref{ar}, onde $\epsilon_t$
representa um termo de erro aleatório:

\begin{equation}
    \label{ar}
    X_t=\alpha_1X_{t-1} + \alpha_2X_{t-2} +(...) + \alpha_pX_{t-p} + \epsilon_t
\end{equation}

Já a componente de média model indica a relação de dependência de uma observação e o erro residual em pontos anteriores.
Isto é, $X_t$ depende de perturbações aleatórias anteriores $\epsilon_{t-i}$ moduladas por coeficientes $\theta_i$,
conforme \autoref{ma}:

\begin{equation}
    \label{ma}
    X_t=\theta_1\epsilon_{t-1} + \theta_2\epsilon_{t-2} + (...) + \theta_q\epsilon_{t-q}
\end{equation}

A combinação dos dois componentes forma o modelo ARMA, como mostrado na equação \autoref{arma}:

\begin{equation}
    \label{arma}
    X_t = \epsilon_t + \sum_{i=1}^{p}\phi_i X_{t-i} + \sum_{i=1}^q \theta_i \epsilon_{t-i}
\end{equation}

Entretanto, modelos ARMA pressupõem que a série seja estacionária (com média e variância constantes ao longo do
tempo). Na prática, muitas séries reais apresentam tendências. Para tratar essa não-estacionariedade, o modelo ARIMA
introduz o parâmetro $d$, referente à ordem de integração.

O parâmetro $d$ indica o número de vezes que a série deve ser diferenciada para tornar-se estacionária. A primeira
diferença ($d=1$) é definida como $\Delta X_t = X_t - X_{t-1}$. Para facilitar a manipulação matemática, utiliza-se o
operador de defasagem (ou \textit{lag}) $L$, definido como $L^i X_t = X_{t-i}$. Assim, a diferenciação de ordem $d$ é
representada pelo operador $(1 - L)^d$.

Ao aplicar este operador à série original, transformamos $X_t$ em uma série estacionária antes de aplicar os componentes
AR e MA. O modelo ARIMA($p, d, q$) resultante é expresso na \autoref{arima}:

\begin{equation}
    \label{arima}
    \left(1 - \sum_{i=1}^p \phi_i L^i\right) (1 - L)^d X_t = \left(1 + \sum_{i=1}^q \theta_i L^i\right) \epsilon_t
\end{equation}

A extensão SARIMA incorpora componentes sazonais adicionais para lidar com padrões que se repetem em intervalos fixos,
como meses ou trimestres. Esses termos sazonais seguem a mesma estrutura do ARIMA, com parâmetros $P$, $D$, $Q$ e
período de sazonalidade $s$. Essa estrutura permite ao modelo capturar tanto as dependências de curto prazo quanto os
padrões cíclicos de longo prazo presentes em séries temporais com sazonalidade definida.

% \begin{equation} \label{sarima} \Phi_P(L^s) \phi_p(L)(1 - L)^d (1 - L^s)^D X_t = \Theta_Q(L^s) \theta_q(L) \epsilon_t
%     \end{equation}

% Nessa expressão $\phi_p(L)$ e $\theta_q(L)$ são os polinômios não sazonais AR e MA, respectivamente, $\Phi_P(L^s)$ e
% $\Theta_Q(L^s)$ são os polinômios sazonais AR e MA, aplicados com defasagem múltipla de $s$ períodos e $(1 - L)^d$
% representa a diferenciação não sazonal e $(1 - L^s)^D$ a diferenciação sazonal.